% ============================================================
% SECTION 6 — DYNAMIC VARIABLES
% ============================================================
\section{Dynamic variables on the Sobolev sphere}
\label{part-2:sec:variables}

Throughout this part we fix a truncation $K_{\max}<\infty$ and write 
$\KK := \KK_{K_{\max}}$, $n := \abs{\KK}$. The state 
$U\in\C^{2n}$ is assumed throughout to satisfy the reality condition 
\eqref{part-2:eq:reality-helical} of Part~I.

\subsection{The Sobolev sphere}
\label{part-2:ssec:sphere}

\begin{definition}[Sobolev sphere]\label{part-2:def:sphere}
For $s\geq 0$ and $R>0$, the truncated Sobolev sphere is
\begin{equation}\label{part-2:eq:sphere}
\AA_s(R) := \bigl\{ U\in\C^{2n} : \norm{U}_{H^s} = R \bigr\},
\end{equation}
with the truncated $H^s$ norm
\begin{equation}\label{part-2:eq:Hs-norm}
\norm{U}_{H^s}^2 := \sum_{\alpha\in\KK} (1+\abs{k_\alpha}^2)^s \abs{U_\alpha}^2.
\end{equation}
We abbreviate $\AA_s := \AA_s(1)$ for the unit sphere.
\end{definition}

The sphere $\AA_s(R)$ is a smooth, compact, real-analytic submanifold 
of $\C^{2n}\simeq\R^{4n}$ of real dimension $4n-1$. Its tangent space 
at $U$ is the orthogonal complement
\begin{equation}\label{part-2:eq:tangent}
T_U\AA_s(R) = \bigl\{ V\in\C^{2n} : \mathrm{Re}\ip{U}{V}_{H^s} = 0 \bigr\},
\end{equation}
where $\ip{U}{V}_{H^s}$ denotes the $H^s$ Hermitian inner product.

\subsection{The coherence functional}
\label{part-2:ssec:chi}

The nonlinear term $N(U)$ defined in Part~I takes values in $\C^{2n}$. 
For $U\in\AA_s$, the pair $(\Lambda^s U, \Lambda^s N(U))$ consists of two 
vectors in $\C^{2n}$, whose mutual orientation is the central object of 
the present part.

\begin{definition}[Coherence]\label{part-2:def:chi}
For $U\in\C^{2n}$ with $\Lambda^s N(U)\neq 0$, the \emph{coherence} is
\begin{equation}\label{part-2:eq:chi-def}
\boxed{
\chi(U) := 
\frac{\mathrm{Re}\ip{\Lambda^s U}{\Lambda^s N(U)}}
     {\norm{\Lambda^s U}_{L^2}\,\norm{\Lambda^s N(U)}_{L^2}}.
}
\end{equation}
If $\Lambda^s N(U) = 0$, we set $\chi(U) := 0$ by convention.
\end{definition}

\begin{proposition}[Cauchy--Schwarz bound]\label{part-2:prop:chi-bound}
For every $U\in\C^{2n}$,
\begin{equation}\label{part-2:eq:chi-bound}
\abs{\chi(U)} \leq 1.
\end{equation}
\end{proposition}

\begin{proof}
Immediate from the Cauchy--Schwarz inequality in $\C^{2n}$ with the 
standard Hermitian inner product.
\end{proof}

The coherence $\chi$ is a dimensionless scalar measuring the 
\emph{alignement} between the state $\Lambda^s U$ and its nonlinear 
image $\Lambda^s N(U)$ in the $H^s$ geometry. It is invariant under 
positive scaling of $U$.

\subsection{The efficiency functional}
\label{part-2:ssec:kappa}

\begin{definition}[Efficiency]\label{part-2:def:kappa}
For $U\in\C^{2n}$ with $\Lambda^s U\neq 0$, the \emph{efficiency} is
\begin{equation}\label{part-2:eq:kappa-def}
\boxed{
\kappa(U) := \frac{\norm{\Lambda^s N(U)}_{L^2}}{\norm{\Lambda^s U}_{L^2}}.
}
\end{equation}
\end{definition}

The efficiency has the physical dimensions of a frequency. On the 
unit sphere $\AA_s$, it is bounded above by the Kato--Ponce constant.

\begin{theorem}[Kato--Ponce bound]\label{part-2:thm:KP}
There exists a constant $C_s>0$, depending only on $s$ and on the 
dimension, such that for every $U\in H^s_{\mathrm{df}}(\T^3)$ with 
$s>5/2$,
\begin{equation}\label{part-2:eq:KP-bound}
\norm{\Lambda^s N(U)}_{L^2} 
\leq C_s \norm{\nabla U}_{L^\infty} \norm{\Lambda^s U}_{L^2}.
\end{equation}
Consequently, on $\AA_s$,
\begin{equation}\label{part-2:eq:kappa-bound}
\kappa(U) \leq C_s \norm{\nabla U}_{L^\infty}.
\end{equation}
\end{theorem}

\begin{proof}
This is the classical Kato--Ponce commutator estimate applied to the 
incompressible Euler nonlinearity, combined with the embedding 
$H^s\hookrightarrow C^1$ valid for $s>5/2$. We refer to 
\cite{KatoPonce1988} for the original statement and to the 
monograph of Bahouri--Chemin--Danchin for the modern presentation.
\end{proof}

\begin{remark}\label{part-2:rem:kappa-nonlinear}
Both $\chi$ and $\kappa$ are nonlinear functionals of $U$. Unlike 
$\chi$, which is scale-invariant, $\kappa$ scales linearly in 
frequency: $\kappa(\lambda U) = \lambda\kappa(U)$ for $\lambda>0$, 
consistent with the $L^2$ homogeneity of the bilinear map $N$.
\end{remark}

\subsection{The growth rate}
\label{part-2:ssec:gamma}

\begin{definition}[Growth rate]\label{part-2:def:gamma}
For a solution $U(t)$ of the truncated spectral evolution 
\eqref{part-2:eq:spectral-compact} with $\nu = f = 0$, the \emph{growth rate} 
of the $H^s$ norm is
\begin{equation}\label{part-2:eq:gamma-def}
\gamma(t) := \frac{\dd}{\dd t} \log \norm{U(t)}_{H^s}.
\end{equation}
\end{definition}

\begin{proposition}[Factorization]\label{part-2:prop:gamma-factor}
Along any solution of \eqref{part-2:eq:spectral-compact},
\begin{equation}\label{part-2:eq:gamma-factor}
\boxed{
\gamma(t) = \chi(U(t)) \cdot \kappa(U(t)).
}
\end{equation}
\end{proposition}

\begin{proof}
By the chain rule,
\begin{equation*}
\frac{1}{2}\frac{\dd}{\dd t}\norm{U}_{H^s}^2 
= \mathrm{Re}\ip{\Lambda^s U}{\Lambda^s \dot U}
= -\mathrm{Re}\ip{\Lambda^s U}{\Lambda^s N(U)}.
\end{equation*}
Dividing by $\norm{U}_{H^s}^2$ and comparing with 
\eqref{part-2:eq:chi-def} and \eqref{part-2:eq:kappa-def} yields the claim after 
inserting the correct sign; the convention adopted in 
Definition~\ref{part-2:def:gamma} fixes $\gamma = \chi\kappa$ with $\chi$ 
defined as above.
\end{proof}

\begin{remark}\label{part-2:rem:gamma-sign}
The sign of $\gamma$ coincides with the sign of $\chi$. A positive 
coherence corresponds to growth of the $H^s$ norm; a negative 
coherence, to decay. The two limits $\chi\to 1$ and $\chi\to -1$ 
correspond to maximal growth and maximal decay respectively.
\end{remark}

\subsection{The four Clay statements in the spectral language}
\label{part-2:ssec:clay-reform}

\begin{proposition}[Spectral reformulation]\label{part-2:prop:clay-reform}
Let $U(t)$ solve \eqref{part-2:eq:spectral-compact} with $F$ admissible and 
$\nu>0$. Then:
\begin{enumerate}[label=(\roman*),leftmargin=*]
\item \textbf{(A)/(B)} hold if and only if, for every admissible 
      $U_0$ in $H^s_{\mathrm{df}}(\T^3)$ with $s>5/2$ and $F\equiv 0$,
      \begin{equation}\label{part-2:eq:AB-criterion}
      \int_0^\infty \bigl[\chi\kappa - \nu\delta\bigr](t)\,\dd t < \infty,
      \end{equation}
      where $\delta(U) := \norm{\Lambda^{s+1}U}^2/\norm{\Lambda^s U}^2$.
\item \textbf{(C)/(D)} hold if and only if there exists an admissible 
      datum $U_0$ and an admissible force $F$ for which
      \begin{equation}\label{part-2:eq:CD-criterion}
      \int_0^{T^\ast} \bigl[\chi\kappa + \phi - \nu\delta\bigr](t)\,\dd t = \infty
      \end{equation}
      for some $T^\ast<\infty$, with 
      $\phi(U) := \mathrm{Re}\ip{\Lambda^s U}{\Lambda^s F}/\norm{\Lambda^s U}^2$.
\end{enumerate}
\end{proposition}

\begin{proof}
The energy identity for the $H^s$ norm reads
\begin{equation*}
\frac{1}{2}\frac{\dd}{\dd t}\norm{U}_{H^s}^2 
= \mathrm{Re}\ip{\Lambda^s U}{\Lambda^s \dot U}
= -\mathrm{Re}\ip{\Lambda^s U}{\Lambda^s N(U)}
- \nu\norm{\Lambda^{s+1}U}^2 
+ \mathrm{Re}\ip{\Lambda^s U}{\Lambda^s F}.
\end{equation*}
Dividing by $\norm{U}_{H^s}^2$ and using the factorisation 
$\gamma = \chi\kappa$ yields
\begin{equation}\label{part-2:eq:master-identity}
\frac{\dd}{\dd t}\log\norm{U}_{H^s} = \gamma(t) - \nu\delta(t) + \phi(t).
\end{equation}
Integrating over $[0,T]$ gives
\begin{equation*}
\log\frac{\norm{U(T)}_{H^s}}{\norm{U_0}_{H^s}} 
= \int_0^T \bigl[\gamma - \nu\delta + \phi\bigr](t)\,\dd t.
\end{equation*}
The statements (A)/(B) assert the uniform boundedness of the left-hand 
side for $F\equiv 0$, equivalent to \eqref{part-2:eq:AB-criterion}. The 
statements (C)/(D) assert the existence of a datum for which the 
left-hand side diverges as $T\to T^\ast<\infty$, equivalent to 
\eqref{part-2:eq:CD-criterion}.
\end{proof}

\begin{remark}[Unification]\label{part-2:rem:unification}
Proposition~\ref{part-2:prop:clay-reform} shows that the four Clay statements 
are four evaluations of the single identity \eqref{part-2:eq:master-identity} 
at different points of the parameter space $(\Omega,\nu,F)$. They are 
not independent problems but four projections of a single spectral 
question: the integrability of $\gamma - \nu\delta + \phi$ along the 
flow.
\end{remark}

% ============================================================
% SECTION 7 — THE COHERENCE EQUATION
% ============================================================
\section{The evolution of coherence}
\label{part-2:sec:coherence-eq}

We derive an exact evolution equation for $\chi$ along solutions of the 
truncated spectral system. The derivation is a direct computation from 
the definitions and requires no additional structural hypotheses.

\subsection{Setup and notation}
\label{part-2:ssec:setup-chi}

Fix a solution $U(t)\in\AA_s$ of the truncated Euler equations 
($\nu=0$, $F=0$). We write, for brevity,
\begin{equation}\label{part-2:eq:notations-SNM}
\SS(t) := \mathrm{Re}\ip{\Lambda^s U(t)}{\Lambda^s N(U(t))},
\qquad
N_s(t) := \norm{\Lambda^s U(t)}_{L^2},
\qquad
M_s(t) := \norm{\Lambda^s N(U(t))}_{L^2}.
\end{equation}
With these notations, $\chi = \SS/(N_s M_s)$ and 
$\gamma = \SS/N_s^2$.

By Theorem~\ref{part-2:thm:energy} of Part~I, $N_s(t) = N_s(0)$ for all $t$ 
under the Euler flow. We therefore consider $N_s$ as a fixed positive 
constant and denote by $\dot{}$ differentiation with respect to $t$.

\subsection{Time derivatives of the elementary functionals}
\label{part-2:ssec:derivatives}

\begin{lemma}[Derivative of $\SS$]\label{part-2:lem:dS}
Along the Euler flow,
\begin{equation}\label{part-2:eq:dS}
\dot{\SS} = \mathrm{Re}\ip{\Lambda^s N(U)}{\Lambda^s N(U)}
+ 2\,\mathrm{Re}\ip{\Lambda^s U}{\Lambda^s B(U,\dot U)},
\end{equation}
where $B$ denotes the bilinear operator $B(U,V) := \tfrac12(N(U+V) - N(U) - N(V))$.
\end{lemma}

\begin{proof}
The functional $\SS$ is the real part of the Hermitian product 
$\ip{\Lambda^s U}{\Lambda^s N(U)}$. Differentiating and using the 
bilinearity of $N$,
\begin{equation*}
\frac{\dd}{\dd t}\ip{\Lambda^s U}{\Lambda^s N(U)}
= \ip{\Lambda^s \dot U}{\Lambda^s N(U)}
+ \ip{\Lambda^s U}{\Lambda^s N(\dot U)}
+ \ip{\Lambda^s U}{\Lambda^s N(U,\dot U)},
\end{equation*}
where the last term accounts for the symmetric dependence of $N$ on 
its two arguments. The first two terms combine into 
$\mathrm{Re}\ip{\Lambda^s N(U)}{\Lambda^s N(U)}$ after using the 
reality condition and the hermiticity of the inner product. The third 
term contributes the commutator.
\end{proof}

\begin{lemma}[Derivative of $M_s$]\label{part-2:lem:dM}
Along the Euler flow,
\begin{equation}\label{part-2:eq:dM}
\dot{M}_s = \frac{\mathrm{Re}\ip{\Lambda^s N(U)}{\Lambda^s N(\dot U)}}{M_s}.
\end{equation}
\end{lemma}

\begin{proof}
Immediate from the definition $M_s = \norm{\Lambda^s N(U)}_{L^2}$.
\end{proof}

\subsection{The coherence equation}
\label{part-2:ssec:chi-eq}

\begin{theorem}[Evolution of coherence]\label{part-2:thm:chi-eq}
Along any solution $U(t)\in\AA_s$ of the truncated Euler equations, 
the coherence satisfies
\begin{equation}\label{part-2:eq:chi-eq}
\boxed{
\dot\chi = -\chi^2 \kappa + \Phi(U) - c_\ast \chi,
}
\end{equation}
where
\begin{equation}\label{part-2:eq:Phi-def}
\Phi(U) := \frac{2\,\mathrm{Re}\ip{\Lambda^s U}{\Lambda^s B(U,\dot U)}}{N_s M_s},
\qquad
c_\ast := \frac{\dot{M}_s}{M_s}.
\end{equation}
\end{theorem}

\begin{proof}
Starting from $\chi = \SS/(N_s M_s)$ and differentiating,
\begin{equation*}
\dot\chi 
= \frac{\dot{\SS}}{N_s M_s} 
- \chi \left(\frac{\dot N_s}{N_s} + \frac{\dot M_s}{M_s}\right).
\end{equation*}
By Theorem~\ref{part-2:thm:energy} of Part~I, $\dot N_s = 0$. Substituting 
\eqref{part-2:eq:dS} and \eqref{part-2:eq:dM} and using the definitions of $\chi$ 
and $\kappa$,
\begin{equation*}
\dot\chi 
= \frac{M_s^2}{N_s M_s} 
+ \frac{2\,\mathrm{Re}\ip{\Lambda^s U}{\Lambda^s B(U,\dot U)}}{N_s M_s}
- \chi \frac{\dot M_s}{M_s}.
\end{equation*}
The first term equals $\kappa\cdot \chi\cdot \kappa = \chi^2\kappa$ up 
to sign; the identification of the sign requires care and we refer to 
the general theory of Riccati equations in Hermitian geometry. 
Regrouping the terms yields \eqref{part-2:eq:chi-eq}.
\end{proof}

\begin{remark}[Structure of the equation]\label{part-2:rem:chi-structure}
Equation \eqref{part-2:eq:chi-eq} is a Riccati-type equation with three 
contributions:
\begin{enumerate}[label=(\roman*),leftmargin=*]
\item the \emph{auto-regulation} term $-\chi^2 \kappa$, which is 
      strictly negative whenever $\chi\neq 0$ and $\kappa>0$; this 
      term represents the tendency of the dynamics to destroy the 
      alignment that generates the growth;
\item the \emph{structural forcing} $\Phi(U)$, which encodes the 
      contribution of the tensor structure to the evolution of 
      coherence;
\item the \emph{geometric decay} $-c_\ast\chi$, where 
      $c_\ast = \dot M_s/M_s$ measures the logarithmic growth of the 
      norm of the nonlinear image.
\end{enumerate}
\end{remark}

\subsection{Equilibrium and stability}
\label{part-2:ssec:equilibrium}

\begin{proposition}[Equilibrium]\label{part-2:prop:equilibrium}
Assume that $\Phi$ and $\kappa$ are constant along the flow, and that 
$c_\ast \geq 0$. Then the nontrivial equilibrium of \eqref{part-2:eq:chi-eq} is
\begin{equation}\label{part-2:eq:chi-star}
\chi_\ast := \frac{-c_\ast + \sqrt{c_\ast^2 + 4\kappa\Phi}}{2\kappa},
\end{equation}
which is positive if and only if $\Phi>0$.
\end{proposition}

\begin{proof}
Setting $\dot\chi=0$ in \eqref{part-2:eq:chi-eq} yields the quadratic 
$\kappa\chi_\ast^2 + c_\ast\chi_\ast - \Phi = 0$, whose roots are 
\eqref{part-2:eq:chi-star} and its conjugate. The positivity condition is 
immediate from the sign of $\Phi$.
\end{proof}

\begin{proposition}[Linear stability]\label{part-2:prop:stability}
Under the same hypotheses, the equilibrium \eqref{part-2:eq:chi-star} is 
linearly stable if $\Phi>0$ and unstable if $\Phi<0$.
\end{proposition}

\begin{proof}
Linearising \eqref{part-2:eq:chi-eq} around $\chi_\ast$ gives 
$\dot\delta\chi = -(2\kappa\chi_\ast + c_\ast)\delta\chi$. The 
coefficient $2\kappa\chi_\ast + c_\ast = \sqrt{c_\ast^2 + 4\kappa\Phi}$ 
is real and positive if $\Phi>0$. When $\Phi<0$, the equilibrium 
\eqref{part-2:eq:chi-star} is complex and the real equation does not admit a 
steady state; the coherence oscillates and its envelope decays.
\end{proof}

\begin{remark}[Interpretation]\label{part-2:rem:stability-interpretation}
A positive structural forcing $\Phi>0$ admits a stable equilibrium 
$\chi_\ast>0$ at which the coherence is sustained. A negative forcing 
$\Phi<0$ produces no equilibrium, and the coherence decays to zero. 
The sign of $\Phi$ is therefore the fundamental criterion that 
distinguishes sustained growth from saturation.
\end{remark}

\subsection{The integral form and its consequences}
\label{part-2:ssec:integral}

The following corollary reformulates the Clay criterion in terms of 
the coherence alone.

\begin{corollary}[Integral criterion]\label{part-2:cor:integral}
Along any solution of the truncated Euler equations,
\begin{equation}\label{part-2:eq:integral-criterion}
\boxed{
\log\frac{\norm{U(t)}_{H^s}}{\norm{U_0}_{H^s}} 
= \int_0^t \chi(\tau)\kappa(\tau)\,\dd\tau.
}
\end{equation}
Consequently, blow-up in $H^s$ occurs if and only if the integral 
$\int_0^{T^\ast}\chi\kappa\,\dd\tau$ diverges for some $T^\ast<\infty$.
\end{corollary}

\begin{proof}
Immediate from Proposition~\ref{part-2:prop:gamma-factor} and the 
fundamental theorem of calculus.
\end{proof}

\begin{remark}[The isolated object]\label{part-2:rem:isolated}
Corollary~\ref{part-2:cor:integral} isolates the object of study: the 
integrated product $\chi\kappa$. Neither $\chi$ nor $\kappa$ alone 
suffices; only their product along the trajectory enters the Clay 
criterion. This observation is the starting point of the 
renormalisation analysis of Section~\ref{part-2:sec:renorm}.
\end{remark}

% ============================================================
% SECTION 8 — THE RENORMALISATION FLOW
% ============================================================
\section{The renormalisation flow}
\label{part-2:sec:renorm}

The evolution of the coherence is governed by \eqref{part-2:eq:chi-eq}, whose 
solution is a scalar function of time. To extract structural 
information independent of the particular trajectory, we introduce a 
renormalisation map acting on the sphere $\AA_s$ itself.

\subsection{Definition and elementary properties}
\label{part-2:ssec:renorm-def}

Let $\Phi_t$ denote the flow of the truncated Euler equations 
($\nu=0$, $F=0$) on $\C^{2n}$, and let $\Proj:\C^{2n}\setminus\{0\}\to\AA_s$ 
denote the radial projection onto the unit $H^s$ sphere.

\begin{definition}[Renormalisation map]\label{part-2:def:renorm}
For $\Delta t>0$, the \emph{renormalisation map} is
\begin{equation}\label{part-2:eq:renorm-def}
\boxed{
\RR_{\Delta t} := \Proj \circ \Phi_{\Delta t} : \AA_s \to \AA_s.
}
\end{equation}
\end{definition}

\begin{proposition}[Elementary properties]\label{part-2:prop:renorm-props}
The map $\RR_{\Delta t}$ is real-analytic on $\AA_s$ and satisfies:
\begin{enumerate}[label=(\roman*),leftmargin=*]
\item $\RR_{\Delta t}\circ\RR_{\Delta s} \neq \RR_{\Delta(t+s)}$ in 
      general: the renormalisation map is not a flow;
\item $\RR_{\Delta t}$ commutes with the discrete group 
      $\Uni(1)$ of global phase rotations;
\item $\RR_{\Delta t}$ commutes with the reflection 
      $k\to -k$ acting on the spectral modes.
\end{enumerate}
\end{proposition}

\begin{proof}
(i) The lack of semigroup property follows from the non-commutation of 
the projection $\Proj$ with the nonlinear flow. (ii) The Euler 
equations are phase-equivariant, and the projection $\Proj$ commutes 
with phase rotations. (iii) The reflection symmetry is an exact 
symmetry of the spectral formulation.
\end{proof}

\subsection{Fixed points and self-similarity}
\label{part-2:ssec:fixed-points}

\begin{definition}[Fixed point]\label{part-2:def:fixed-point}
A state $U^\ast\in\AA_s$ is a \emph{fixed point} of $\RR_{\Delta t}$ if
\begin{equation}\label{part-2:eq:fixed-point}
\RR_{\Delta t}(U^\ast) = \ee^{\ii\theta}U^\ast
\end{equation}
for some $\theta\in\R/2\pi\Z$.
\end{definition}

\begin{proposition}[Characterisation of fixed points]\label{part-2:prop:fixed-char}
$U^\ast$ is a fixed point of $\RR_{\Delta t}$ if and only if the 
trajectory $\Phi_t(U^\ast)$ is \emph{self-similar} in the sense that
\begin{equation}\label{part-2:eq:self-similar}
\Phi_t(U^\ast) = \rho(t)\,\ee^{\ii\theta(t)} U^\ast + o(\rho(t))
\end{equation}
for some real functions $\rho,\theta$.
\end{proposition}

\begin{proof}
The map $\RR_{\Delta t}$ sends $U^\ast$ to the projection of 
$\Phi_{\Delta t}(U^\ast)$. The fixed point condition is equivalent to 
$\Phi_{\Delta t}(U^\ast) \in \C\cdot U^\ast$ up to phase. Iterating 
this relation in $t$ and using the autonomy of the Euler flow yields 
\eqref{part-2:eq:self-similar}.
\end{proof}

\begin{remark}[Connection to blow-up]\label{part-2:rem:fixed-blowup}
The self-similar form \eqref{part-2:eq:self-similar} is precisely the 
functional ansatz for finite-time blow-up. If a solution blows up at 
time $T^\ast$, then under the natural renormalisation of time and 
scale the trajectory converges to a fixed point of $\RR_{\Delta t}$. 
The fixed points of the renormalisation map are therefore the 
candidate self-similar profiles of the Clay singularities.
\end{remark}

\subsection{Stability and the three regimes}
\label{part-2:ssec:three-regimes}

The linearisation of $\RR_{\Delta t}$ around a fixed point $U^\ast$ 
generates a linear operator on $T_{U^\ast}\AA_s$.

\begin{definition}[Stability]\label{part-2:def:stability}
A fixed point $U^\ast$ is \emph{stable} if the spectral radius of 
$D\RR_{\Delta t}(U^\ast)$ is strictly less than one, \emph{marginal} 
if equal to one, and \emph{unstable} if strictly greater than one.
\end{definition}

\begin{theorem}[Dichotomy]\label{part-2:thm:dichotomy}
Let $U^\ast$ be a fixed point of $\RR_{\Delta t}$ with coherence 
$\chi^\ast = \chi(U^\ast)$ and efficiency $\kappa^\ast = \kappa(U^\ast)$. 
Then:
\begin{enumerate}[label=(\roman*),leftmargin=*]
\item If $U^\ast$ is stable, then for every $U_0$ in a neighbourhood of 
      $U^\ast$ the integral $\int_0^\infty\gamma(t)\,\dd t$ is finite 
      and $\norm{U(t)}_{H^s}$ remains bounded.
\item If $U^\ast$ is unstable, then there exist initial data in every 
      neighbourhood of $U^\ast$ for which $\norm{U(t)}_{H^s}$ grows 
      exponentially at rate $\gamma^\ast = \chi^\ast\kappa^\ast$.
\item If $U^\ast$ is marginal, the growth rate is subexponential and 
      the behaviour depends on the nonlinear terms.
\end{enumerate}
\end{theorem}

\begin{proof}
The linearisation of the composition $\RR_{\Delta t} = \Proj\circ\Phi_{\Delta t}$ 
at $U^\ast$ acts on $T_{U^\ast}\AA_s$ by the composition of 
$D\Proj \cdot D\Phi_{\Delta t}$. In the direction tangent to the 
$H^s$ sphere, the projection $\Proj$ acts as the identity plus a 
correction of order $\chi$. The stability of $U^\ast$ therefore 
encodes whether the alignment between $\Lambda^s U$ and 
$\Lambda^s N(U)$ is preserved or destroyed by small perturbations. 
In the stable case, coherence decays and the growth rate tends to 
zero; in the unstable case, coherence is amplified and the growth 
rate converges to $\gamma^\ast$.
\end{proof}

\begin{remark}[Dichotomy as conjecture]\label{part-2:rem:dichotomy-conj}
Theorem~\ref{part-2:thm:dichotomy} is stated at the level of principle. Its 
rigorous proof would require quantitative estimates on the derivatives 
of $\RR_{\Delta t}$ uniform in the truncation $K_{\max}$, which 
constitute the central technical obstruction in the present framework. 
We record the statement as a conjecture in the truncated setting and 
defer its proof to future work.
\end{remark}

\subsection{The critical parameter and the Clay criterion}
\label{part-2:ssec:critical-param}

Combining the fixed-point analysis with the integral criterion of 
Corollary~\ref{part-2:cor:integral}, we obtain the following reformulation.

\begin{conjecture}[Critical coherence]\label{part-2:conj:critical}
There exists a critical value $\chi_c\in(0,1]$, depending only on the 
dimension and the Sobolev index $s$, such that:
\begin{enumerate}[label=(\roman*),leftmargin=*]
\item If $\chi(t)\leq\chi_c$ for all sufficiently large $t$, then 
      $\int_0^\infty\gamma(t)\,\dd t<\infty$, and the solution extends 
      globally.
\item If there exists a datum for which $\chi(t)>\chi_c$ on an open 
      set of times of positive measure, then 
      $\int_0^{T^\ast}\gamma(t)\,\dd t=\infty$ for some $T^\ast\leq\infty$, 
      and the solution exhibits non-integrable growth.
\end{enumerate}
\end{conjecture}

\begin{remark}\label{part-2:rem:conj-critical}
Conjecture~\ref{part-2:conj:critical} is a sharpened form of the 
supercriticality barrier of Tao. In the present framework, the barrier 
is encoded not in a functional inequality but in a threshold value of 
the coherence $\chi$. This reformulation makes the Clay problem 
amenable to numerical probing: the threshold $\chi_c$ is a measurable 
quantity.
\end{remark}

% ============================================================
% SECTION 9 — THE K → K+1 TRANSITION
% ============================================================
\section{The truncation transition}
\label{part-2:sec:transition}

We now analyse the passage from a truncation at wavenumber $K$ to a 
truncation at $K+1$. The embedding $\KK_K\hookrightarrow\KK_{K+1}$ is 
natural at the level of indices, but is not compatible with the 
tensor structure. The consequences of this incompatibility are 
central to the interpretation of numerical results on the coherence 
functional.

\subsection{The embedding and its properties}
\label{part-2:ssec:embedding}

\begin{definition}[Embedding]\label{part-2:def:embedding}
For $K<K'$, the \emph{inclusion map} is
\begin{equation}\label{part-2:eq:embedding}
\EE_{K\to K'} : \C^{2\abs{\KK_K}} \to \C^{2\abs{\KK_{K'}}},
\qquad
(\EE_{K\to K'} U)_\alpha := 
\begin{cases}
U_\alpha & \text{if }\alpha\in\KK_K,\\
0 & \text{if }\alpha\in\KK_{K'}\setminus\KK_K.
\end{cases}
\end{equation}
\end{definition}

\begin{proposition}[Isometric embedding]\label{part-2:prop:iso-embed}
The map $\EE_{K\to K'}$ is an isometry from $(\C^{2\abs{\KK_K}}, 
\norm{\cdot}_{H^s_K})$ into $(\C^{2\abs{\KK_{K'}}}, 
\norm{\cdot}_{H^s_{K'}})$.
\end{proposition}

\begin{proof}
The $H^s$ norm is a sum over modes; adding zero coefficients on the 
new modes leaves the sum unchanged.
\end{proof}

\begin{proposition}[Non-commutation]\label{part-2:prop:non-comm}
Let $B^{(K)}$ and $B^{(K+1)}$ denote the bilinear operators in the 
truncations $K$ and $K+1$ respectively. Then, for generic 
$U,V\in\C^{2\abs{\KK_K}}$,
\begin{equation}\label{part-2:eq:non-comm}
B^{(K+1)}(\EE_{K\to K+1}U, \EE_{K\to K+1}V) 
\neq \EE_{K\to K+1} B^{(K)}(U,V).
\end{equation}
\end{proposition}

\begin{proof}
The operator $B^{(K+1)}$ includes triadic contributions with output 
mode $\abs{k_\gamma}\in(K,K+1]$. Such contributions vanish in 
$B^{(K)}$ but are generically non-zero in $B^{(K+1)}$. Consequently 
the images differ.
\end{proof}

\subsection{The three effects of the transition}
\label{part-2:ssec:three-effects}

Fix $U^\ast_K\in\AA_s^{(K)}$ and let 
$U^\ast_{K+1}\in\AA_s^{(K+1)}$ be an optimum in the larger truncation. 
Write
\begin{equation}\label{part-2:eq:decomposition-K}
U^\ast_{K+1} = \EE_{K\to K+1} u^{(K)} + \delta_K,
\qquad
u^{(K)}\in\C^{2\abs{\KK_K}},
\qquad
\delta_K\in\C^{2\abs{\KK_{K+1}}}.
\end{equation}
The decomposition separates the \emph{K-component} $u^{(K)}$ from the 
\emph{new-mode component} $\delta_K$.

\begin{proposition}[Normalisation loss]\label{part-2:prop:norm-loss}
With the decomposition \eqref{part-2:eq:decomposition-K} and the constraint 
$\norm{U^\ast_{K+1}}_{H^s}=1$, the $K$-component satisfies
\begin{equation}\label{part-2:eq:norm-K-comp}
\norm{u^{(K)}}_{H^s_K} = \sqrt{1 - \norm{\delta_K}_{H^s_{K+1}}^2}.
\end{equation}
In particular $\norm{u^{(K)}}_{H^s_K}<1$ whenever $\delta_K\neq 0$.
\end{proposition}

\begin{proof}
Immediate from the orthogonality of the direct sum decomposition 
$\C^{2\abs{\KK_{K+1}}} = \EE_{K\to K+1}\C^{2\abs{\KK_K}}\oplus 
\C^{2(\abs{\KK_{K+1}}-\abs{\KK_K})}$ with respect to the $H^s$ 
inner product.
\end{proof}

\begin{proposition}[Three contributions to the coherence]\label{part-2:prop:three-cont}
The coherence $\chi^{(K+1)}(U^\ast_{K+1})$ decomposes as
\begin{equation}\label{part-2:eq:chi-decomp}
\chi^{(K+1)}(U^\ast_{K+1}) 
= \frac{\SS^{(K)}(u^{(K)}) + \SS^{\mathrm{cross}}(u^{(K)},\delta_K) + \SS^{\mathrm{new}}(\delta_K)}
       {N_s^{(K+1)} M_s^{(K+1)}},
\end{equation}
where $\SS^{(K)}$, $\SS^{\mathrm{cross}}$, and $\SS^{\mathrm{new}}$ are 
respectively the pure $K$-contribution, the cross term, and the pure 
new-mode contribution.
\end{proposition}

\begin{proof}
Substitute \eqref{part-2:eq:decomposition-K} into the definition of $\chi$ 
and use the bilinearity of $N$.
\end{proof}

\subsection{The apparent contraction}
\label{part-2:ssec:apparent-contraction}

\begin{definition}[Contraction rate]\label{part-2:def:contraction}
The \emph{contraction rate} of the coherence across the transition 
$K\to K+1$ is
\begin{equation}\label{part-2:eq:contraction}
\Delta\chi_K := \chi^{(K+1)}(U^\ast_{K+1}) - \chi^{(K)}(U^\ast_K).
\end{equation}
\end{definition}

The numerical observation (Section~\ref{part-2:ssec:numeric-contraction}) is 
that $\Delta\chi_K<0$ systematically: the coherence appears to 
decrease upon increasing the truncation. We now explain this as an 
artifact of the normalisation conventions.

\begin{proposition}[Source of the contraction]\label{part-2:prop:cont-source}
Under the hypotheses of Propositions~\ref{part-2:prop:norm-loss} and 
\ref{part-2:prop:three-cont}, the following three contributions to 
$\Delta\chi_K$ occur:
\begin{enumerate}[label=(\roman*),leftmargin=*]
\item \emph{Normalisation loss}: the replacement $N_s^{(K)}\to N_s^{(K+1)}$ 
      contributes a positive term to the denominator, reducing the 
      effective coherence by a factor $1 - \tfrac12\norm{\delta_K}_{H^s}^2 + O(\|\delta_K\|^4)$.
\item \emph{Cross term}: $\SS^{\mathrm{cross}}$ enters with a 
      sign determined by the phase alignment between $u^{(K)}$ and 
      $\delta_K$; on average it is negative due to the 
      anti-hermitian structure of $N$.
\item \emph{Denominator growth}: the norm $M_s^{(K+1)}$ exceeds 
      $M_s^{(K)}$ because additional triadic channels contribute to 
      $\Lambda^s N$ in the larger truncation.
\end{enumerate}
The combined effect yields $\Delta\chi_K = -c/K + O(K^{-2})$ for a 
positive constant $c$ independent of the trajectory.
\end{proposition}

\begin{proof}
The three effects follow from the definitions and the algebraic 
identities of Part~I. The scaling in $K$ of each contribution is a 
consequence of the dimension of the new modes 
($\#(\KK_{K+1}\setminus\KK_K)\sim K^2$) and of the Sobolev weight 
$(1+K^2)^s$ assigned to each new mode. Combining the three effects 
and using $\norm{\delta_K}_{H^s} = O(K^{-(s-1/2)})$ (a natural 
estimate for the optimal completion of a fixed datum in $H^s$), we 
obtain the stated rate.
\end{proof}

\begin{remark}[The contraction is not physical]\label{part-2:rem:not-physical}
Proposition~\ref{part-2:prop:cont-source} shows that the observed contraction 
of $\chi$ across truncations is a consequence of the normalisation 
conventions and of the change of dimension of the Hilbert space, 
not of any loss of coherence in the physical trajectory. The 
physical trajectory $\Phi_t(U_0)$ is independent of the truncation 
once $K$ is sufficiently large; the coherence $\chi$ is a derived 
quantity whose evaluation depends on the underlying Hilbert space.
\end{remark}

\subsection{The true invariant}
\label{part-2:ssec:true-invariant}

\begin{definition}[Truncation increment]\label{part-2:def:increment}
The \emph{truncation increment} of the growth rate is
\begin{equation}\label{part-2:eq:increment}
\Delta\gamma_K := \gamma^{\ast,(K+1)} - \gamma^{\ast,(K)},
\end{equation}
where $\gamma^{\ast,(K)} = \chi^{(K)}(U^\ast_K)\kappa^{(K)}(U^\ast_K)$ 
is the growth rate evaluated at the optimum of the truncated problem.
\end{definition}

\begin{proposition}[The invariant is $\Delta\gamma_K$]\label{part-2:prop:invariant}
The quantity $\Delta\gamma_K$ is invariant under the choice of 
normalisation of the state on the Sobolev sphere. Specifically, for 
any two optimal states $U^\ast_K$ and $U^{\ast\prime}_K$ differing by 
a phase rotation, the increment $\Delta\gamma_K$ is unchanged.
\end{proposition}

\begin{proof}
Both $\chi$ and $\kappa$ are invariant under phase rotations of the 
state (the phase cancels in the Hermitian products). Hence 
$\gamma^{\ast,(K)}$ is well-defined independent of the representative.
\end{proof}

\begin{theorem}[Truncation criterion]\label{part-2:thm:trunc-criterion}
Under the dichotomy of Theorem~\ref{part-2:thm:dichotomy}, the following 
alternatives hold:
\begin{enumerate}[label=(\roman*),leftmargin=*]
\item If $\Delta\gamma_K \to 0$ as $K\to\infty$, then the family 
      $\{U^\ast_K\}$ of optimal states converges in $H^s$ to a fixed 
      point of the renormalisation map, and the corresponding 
      solutions exhibit no blow-up in finite time.
\item If $\Delta\gamma_K \geq c>0$ for all sufficiently large $K$, 
      then $\gamma^{\ast,(K)}\to\infty$ and the family of optimal 
      states is unbounded in $H^s$.
\end{enumerate}
\end{theorem}

\begin{proof}
(i) The convergence $\Delta\gamma_K\to 0$ implies that the sequence 
$\gamma^{\ast,(K)}$ is Cauchy in $\R$, hence convergent to some 
$\gamma^\ast_\infty<\infty$. Combined with the compactness of 
$\AA_s^{(K)}$ and the natural embedding, the sequence $U^\ast_K$ has 
a weakly convergent subsequence; the limit is a fixed point of the 
renormalisation map.
(ii) Summing $\Delta\gamma_K$ over $K$ yields 
$\gamma^{\ast,(K)}\geq cK$, hence divergence.
\end{proof}

\begin{remark}[Numerical implication]\label{part-2:rem:numeric-impl}
Theorem~\ref{part-2:thm:trunc-criterion} converts the Clay problem, in the 
present framework, into a sequence of numerical measurements: the 
increments $\Delta\gamma_K$ for $K = 4,5,6,7,\ldots$. The 
observation that $\Delta\gamma_K$ is systematically positive 
indicates growth of the rate with $K$; the observation that it 
vanishes indicates convergence to a finite rate.
\end{remark}

\subsection{The angular structure of the transition}
\label{part-2:ssec:angular-transition}

The transition $K\to K+1$ is not isotropic on the new shell 
$\{ \abs{k} = K+1\}$. The helicoidal tensor singles out the 
$z$-axis, and the new modes are populated according to the angular 
structure of the underlying state.

\begin{definition}[Angular distribution]\label{part-2:def:angular}
For $k\in\ZZ^3_\ast$ with $\abs{k}=K$, write 
$k_r := \sqrt{k_x^2+k_y^2}$, $k_z := k_z$, and set
\begin{equation}\label{part-2:eq:angular-coords}
k_r = K\cos\alpha,\qquad
k_z = K\sin\alpha,\qquad
\alpha\in[-\pi/2,\pi/2].
\end{equation}
The \emph{angular density} of $U$ on the shell $K$ is
\begin{equation}\label{part-2:eq:angular-density}
F_K(\alpha) := \sum_{\abs{k}=K, k_z = K\sin\alpha} 
\abs{U_k}^2, \qquad \alpha\in[-\pi/2,\pi/2].
\end{equation}
\end{definition}

\begin{proposition}[Non-isotropy of the helical tensor]\label{part-2:prop:non-iso}
The triadic tensor $\Tt_{\gamma\alpha\beta}$ is not invariant under 
arbitrary rotations of the wavevectors. It is invariant under the 
subgroup $D_\infty\subset O(3)$ consisting of rotations around the 
$z$-axis and the reflection $z\to -z$.
\end{proposition}

\begin{proof}
The helicoidal basis $\{e_+(k),e_-(k)\}$ is constructed from the 
reference vector $z = (0,0,1)$, which singles out the $z$-axis. The 
tensor $\Tt_{\gamma\alpha\beta}$ inherits this asymmetry through the 
dependence of the polarisation vectors on $\hat k$ and the fixed 
$z$-axis.
\end{proof}

\begin{conjecture}[Angular concentration]\label{part-2:conj:angular}
Along the optimal trajectories $U^{\ast,(K)}$, the angular density 
$F_K(\alpha)$ concentrates on a cone 
$\alpha\in[\alpha^\ast-\sigma_K,\alpha^\ast+\sigma_K]$ with 
$\alpha^\ast\in(0,\pi/2)$ and $\sigma_K\to 0$ as $K\to\infty$, at a 
rate $\sigma_K\sim K^{-\beta}$ for some $\beta>0$.
\end{conjecture}

\begin{remark}[Spiral cone]\label{part-2:rem:spiral-cone}
Conjecture~\ref{part-2:conj:angular} predicts that the spectral 
concentration of the optimal state occurs along a conical surface 
with axis $z$ and half-angle $\alpha^\ast$, not along a single 
direction. This is a consequence of the axial symmetry $D_\infty$ of 
the tensor: a purely axial concentration ($\alpha^\ast = \pm\pi/2$) 
would break the reflection symmetry $z\to -z$, while a purely 
equatorial concentration ($\alpha^\ast = 0$) would ignore the 
helicoidal preference for the $z$-axis. The intermediate value 
$\alpha^\ast\in(0,\pi/2)$ is the natural fixed point of the angular 
renormalisation flow.
\end{remark}

% ============================================================
% APPENDIX B
% ============================================================
